\exercisegroup

\begin{enumerate}
\def\labelenumi{\arabic{enumi})}
\tightlist
\item
  Let B be an orthonormal basis in an infinite dimensional hilbertian space E.
\end{enumerate}

\begin{enumerate}
\def\labelenumi{\alph{enumi})}
\tightlist
\item
  Show that every everywhere dense subset in E has a cardinality at least equal to that of B, and that there exists an everywhere dense set in E which is equipotent to B.
\item
  Show that \(\text{Card}(E) = \text{Card}(B^{N})\) (to see that \(\text{Card}(E) \leqslant \text{Card}(B^{N})\) use a)).
\item
  Show that if \(\text{Card}(B) \leqslant \text{Card}(R)\) , then the cardinality of every algebraic basis of E is equal to \(\text{Card}(R) = 2^{\aleph_{0}}\) (use II, p.~80, exerc. 24, c)); however, if \(\text{Card}(B) > \text{Card}(R)\) , then every algebraic basis of E is equipotent to \(B^{N}\) (use b) and A, II, § 7, exerc. 3, d)).
\end{enumerate}

¶ 2) a) Let \(\mathrm{E}_1, \mathrm{E}_2\) be two hilbertian spaces whose respective hilbertian dimensions are two cardinals \(m\) and \(n\) such that \(m < n \leqslant m^{\aleph_{0}}\) . Let \(\mathrm{E} = \mathrm{E}_1 \oplus \mathrm{E}_2\) be the hilbertian sum of \(\mathrm{E}_1, \mathrm{E}_2\) and let \((b_\lambda)_{\lambda \in L}\) be an orthonormal basis of \(\mathrm{E}_2\) . Show that there exists an algebraically independent system \((a_\lambda)_{\lambda \in L}\) in \(\mathrm{E}_1\) (cf.~exerc. 1, c)); let \(H\) be the subspace of \(E\) generated (algebraically) by the family \((a_\lambda + b_\lambda)_{\lambda \in L}\) . Show that the hilbertian dimension of \(\overline{H}\) is equal to \(H\) (observe that the orthogonal projection from \(H\) onto \(\mathrm{E}_2\) is everywhere dense, and use exerc. 1, a)). If \(S\) is an orthonormal subset of \(H\) , show that \(S \cap E_2 = \varnothing\) ; and deduce that \(\operatorname{Card}(S) \leqslant m\) (observe that every element of an orthonormal basis of \(E_1\) is orthogonal to every element of \(S\) except at most to a countably infinite subset).

\begin{enumerate}
\def\labelenumi{\alph{enumi})}
\setcounter{enumi}{1}
\tightlist
\item
  Let \(E_{3}\) be a hilbertian space with hilbertian dimension \(p \geqslant n\) , and let F be the hilbertian sum \(E \oplus E_{3}\) . Let G be the subspace \(H + E_{3}\) of F. Show that the hilbertian dimension of \(\overline{G}\) is p.~If T is an orthonormal subset of G, show that \(T \cap (E_{2} + E_{3}) \subset E_{3}\) ; deduce that the cardinality of the orthogonal projection from T onto \(E_{2}\) is at most m (argue as in a)). Conclude from this that G does not have an orthonormal basis, by observing that the orthogonal projection from G onto \(E_{2}\) is everywhere dense in \(E_{2}\) .
\end{enumerate}

\begin{enumerate}
\def\labelenumi{\arabic{enumi})}
\setcounter{enumi}{2}
\item
  Show that, in every Hausdorff and non complete prehilbertian space E, there exists a closed hyperplane whose orthogonal subspace in E reduces to 0. Deduce that if E satisfies the first axiom of countability, then there exists a non total orthonormal family in E which is not contained in any orthonormal basis.
\item
  Let E be a Hausdorff prehilbertian space, \((\mathrm{E}_{i})_{i\in\mathbf{I}}\) a family of complete vector subspaces of E, well-ordered by inclusion, such that the union of all \(E_{i}\) is everywhere dense in E. Show that there exists an orthonormal basis \((e_{\alpha})_{\alpha\in\mathbf{A}}\) in E with the following property: for every \(i\in I\) , the set of all \(e_{\alpha}\) belonging to \(E_{i}\) is an orthonormal basis of \(E_{i}\) . (Consider the set of all orthonormal subsets S in E such that, for every \(i\in I\) , every vector of S not belonging to \(E_{i}\) , is orthogonal to \(E_{i}\) , and take a maximal element of this set.) From this, deduce a new proof of the corollary of V, p.~24.
\item
  Show that for a hilbertian space E with an infinite hilbertian dimension, there exists an isomorphism from E onto a closed vector subspace of E which is distinct from E.
\item
  Let E be a hilbertian space and \((e_{i})_{i\in I}\) an orthonormal basis of E. Show that if \((a_{i})_{i\in I}\) is a topologically independent family in E such that \(\sum_{i\in I}\|e_{i}-a_{i}\|^{2}<+\infty\) , then the family \((a_{i})\) is total. (Let J be a finite subset of I; show that there is a continuous linear mapping u from E into itself such that \(u(e_{i})=e_{i}\) for \(i\in J\) , \(u(e_{i})=a_{i}\) for \(i\notin J\) , and that the norm of \(u-1_{E}\) can be made arbitrarily small by a suitable choice of J; then use IV, p.~65, exerc. 17.)
\item
  Given \(n\) points \(x_{i} (1 \leqslant i \leqslant n)\) in a Hausdorff prehilbertian space \(E\) , we mean by the Gram's determinant of these \(n\) points the determinant
\end{enumerate}

\[
\mathrm{G} (x _ {1}, \dots , x _ {n}) = \det (\langle x _ {i} | x _ {j} \rangle).
\]

\begin{enumerate}
\def\labelenumi{\alph{enumi})}
\tightlist
\item
  Show that \(\mathrm{G}(x_{1},\ldots,x_{n})\geqslant0\) and that for \(x_{1},\ldots,x_{n}\) to form an independent system, it is necessary and sufficient that \(\mathrm{G}(x_{1},\ldots,x_{n})\neq0\) (assuming that \(\dim(\mathrm{E})\geqslant n\) , consider an orthogonal basis of an n-dimensional subspace containing \(x_{1},\ldots,x_{n}\) ).
\item
  Show that if \(x_{1},\ldots,x_{n}\) is an independent system in E, the distance from a point \(x\in E\) to the vector subspace V generated by \(x_{1},\ldots,x_{n}\) is equal to \((\mathrm{G}(x,x_{1},\ldots,x_{n})/\mathrm{G}(x_{1},\ldots,x_{n}))^{1/2}\) (find the expression for the orthogonal projection of x on V).
\item
  Let \((x_{n})\) be an infinite sequence of points in E. For the family \((x_{n})\) to be topologically independent, it is necessary and sufficient that, for every integer p\textgreater0,
\end{enumerate}

\[
\sup _ {n} \left(\mathrm{G} \left(x _ {1}, \dots , x _ {p - 1}, x _ {p + 1}, \dots , x _ {n}\right) / \mathrm{G} \left(x _ {1}, \dots , x _ {n}\right)\right) <   + \infty
\]

(use b)).

\begin{enumerate}
\def\labelenumi{\arabic{enumi})}
\setcounter{enumi}{7}
\item
  Let \(\mathbf{E}\) be a hilbertian space which has a countably infinite orthonormal basis \((e_n)_{n\geq 1}\) . Let \(\mathbf{A}\) be the closed convex envelope in \(\mathbf{E}\) of the set consisting of the points \(\left(1 - \frac{1}{n}\right)e_n\) for all \(n \geqslant 1\) . Show that there does not exist any pair of points \(x, y\) in \(A\) whose distance is equal to the diameter of \(A\) (compare with IV, p.~54, exerc. 12).
\item
  \begin{enumerate}
  \def\labelenumii{\alph{enumii})}
  \tightlist
  \item
    Let E be an infinite dimensional real hilbertian space satisfying the first axiom of countability. Let \((a_{n})_{n\geqslant0}\) be a free family of points in E, such that each of the two families \((a_{2n})\) and \((a_{2n+1})\) is total in E (II, p.~80, exerc. 26, a)). Let F and G be two vector subspaces of E for which \((a_{2n})\) and \((a_{2n+1})\) are respectively (algebraic) bases. The spaces F and G are put in separating duality by the bilinear form \(\langle y|z\rangle\) . Show that if B denotes the unit ball in E, then in the space F, endowed with the topology \(\sigma(\mathrm{F},\dot{\mathrm{G}})\) , the convex set \(F\cap B\) is closed, but does not have any closed support hyperplane.
  \end{enumerate}
\end{enumerate}

\begin{enumerate}
\def\labelenumi{\alph{enumi})}
\setcounter{enumi}{1}
\tightlist
\item
  Let \((b_n)_{n\geqslant 1}\) be an everywhere dense set in \(\mathbf{B}\) , and for every \(x\in \mathrm{E}\) , let \(u(x)\) be the sequence \((\langle b_k|x\rangle /k)_{k\geqslant 1}\) . Show that \(u\) is an injective, continuous linear mapping from \(\mathrm{E}\) into the hilbertian space \(\ell_{\mathbf{R}}^2 (\mathbf{N})\) and that \(u(\mathbf{B})\) is compact. Show that, in the normed subspace \(\mathrm{L} = u(\mathrm{F})\) of \(\ell_{\mathbf{R}}^2 (\mathbf{N})\) , the set \(u(\mathbf{B}\cap \mathbf{F})\) is closed, convex and precompact, but does not have any closed support hyperplane (observe that if \(f\) is a continuous linear form on \(\mathbf{L}\) , then \(f\circ u\) is a continuous linear form on \(\mathrm{F}\) for the topology \(\sigma (\mathrm{F},\mathrm{G})\) ).
\end{enumerate}

\begin{enumerate}
\def\labelenumi{\arabic{enumi})}
\setcounter{enumi}{9}
\tightlist
\item
  Let \(\mathbf{E}\) be an infinite dimensional real hilbertian space satisfying the first axiom of countability, and \((e_n)_{n\geqslant 1}\) an orthonormal basis of \(\mathbf{E}\) .
\end{enumerate}

\begin{enumerate}
\def\labelenumi{\alph{enumi})}
\item
  Let \(\mathbf{A}\) be the closed convex balanced envelope of the set of all points \(e_n / n\) in \(\mathbf{E}\) . Show that \(\mathbf{A}\) is compact and that there does not exist any closed supporting hyperplane of \(\mathbf{A}\) at the point 0, but that there exist lines \(\mathbf{D}\) passing through 0 such that \(\mathbf{D} \cap \mathbf{A} = \{0\}\) .
\item
  Let \(\mathbf{F}\) be the hilbertian sum \(\mathbf{E} \oplus \mathbf{R}\) , \(e_0\) a vector which, with the \(e_n\) (for \(n \geqslant 1\) ) forms an orthonormal basis of \(\mathbf{F}\) . If \(\mathbf{B}\) is the closed convex envelope of \(\{e_0\} \cup \mathbf{A}\) , show that there exists a closed segment \(\mathbf{L}\) with mid-point 0 in \(\mathbf{F}\) , such that \(\mathbf{L} \cap \mathbf{B} = \{0\}\) , but there does not exist any closed hyperplane passing through 0 which separates \(\mathbf{L}\) and \(\mathbf{B}\) (even though there exists a closed supporting hyperplane of \(\mathbf{B}\) at 0).
\end{enumerate}

\begin{enumerate}
\def\labelenumi{\arabic{enumi})}
\setcounter{enumi}{10}
\tightlist
\item
  Let \(\mathbf{E}_1, \mathbf{E}_2\) be two infinite dimensional real hilbertian spaces satisfying the first axiom of countability, and E the hilbertian sum \(E_{1} \oplus E_{2}\) (which we shall identify with the product \(E_{1} \times E_{2}\) ). Let \((e_{n})_{n \geqslant 1}\) be an orthonormal basis of \(E_{1}\) ; in \(E_{2}\) , let A be a compact convex set containing 0, and D a line passing through 0, such that \(D \cap A = \{0\}\) and suppose there exists no closed supporting hyperplane of A at the point 0 (exerc. 10). Let \((\alpha_{n})\) , \((\beta_{n})\) be two sequences of numbers \(\geqslant 0\) such that \(\lim_{n \to \infty} \beta_{n} = 0\) and \(\sum_{n} \alpha_{n}^{-1} < 1\) . Let P be the set of all points
\end{enumerate}

\(\sum_{n}\xi_{n}e_{n}\) of \(\mathrm{E}_1\) such that \(0\leqslant \xi_n\leqslant \alpha_n\) for every \(n\geqslant 1\) . Finally, let Q be the closed convex envelope

in E of the set of all points \((\alpha_{n}e_{n}, x + \beta_{n}a)\) , where \(n \geqslant 1\) , \(a \neq 0\) is a fixed point in D and \(x\) ranges over A.

\begin{enumerate}
\def\labelenumi{\alph{enumi})}
\item
  Show that \(\mathbf{P} \cap \mathbf{Q} = \emptyset\) and that there exists no closed hyperplane in E separating \(\mathbf{P}\) and \(\mathbf{Q}\) .
\item
  Let F be the hilbertian sum \(E \oplus R\) , and let c be an arbitrary point of F not contained in E. Show that the pointed convex cones \(P_{1}\) , \(Q_{1}\) with vertex c, generated by P and Q respectively, are closed in F and that there exists no closed hyperplane in F separating \(P_{1}\) and \(Q_{1}\) (to see that \(P_{1}\) and \(Q_{1}\) are closed, prove that neither P nor Q contain the half-line).
\end{enumerate}

\begin{enumerate}
\def\labelenumi{\arabic{enumi})}
\setcounter{enumi}{11}
\item
  Let E be an infinite dimensional real hilbertian space. Show that there exist infinitely many complex hilbertian space structures on E for which E is the real locally convex space underlying these complex hilbertian spaces (II, p.~61). (To prove the existence of the automorphisms u of the topological vector space structure of E, such that \(u^{2}(x) = -x\) , use an orthonormal basis of E; then apply V, p.~60, exerc. 1.) Give an example to show that the proposition does not extend to Hausdorff non complete prehilbertian spaces (consider an everywhere dense hyperplane in such a space).
\item
  Let E be an infinite dimensional hilbertian space satisfying the first axiom of countability and \((e_{n})_{n\in\mathbf{Z}}\) an orthonormal basis of E whose set of indices is the set of rational integers. Let u denote the isometry from E onto itself such that \(u(e_{n}) = e_{n+1}\) for all \(n \in Z\) , and put
\end{enumerate}

\[
f (x) = \frac {1}{2} (1 - \| x \|) e _ {0} + u (x).
\]

\begin{enumerate}
\def\labelenumi{\alph{enumi})}
\item
  Let B be the unit ball and S the unit sphere in E. Show that the restriction of \(f\) to B is a homeomorphism from B onto itself (observe that the restriction of \(u\) to S is a homeomorphism from S onto itself), and that there does not exist any point \(x_0 \in \mathbf{B}\) such that \(f(x_0) = x_0\) (express \(x_0\) in terms of its coordinates with respect to \((e_n)\) ).
\item
  For every \(x \in \mathbf{B}\) , let \(g(x)\) be the point of intersection of \(S\) with the half-line with origin \(f(x)\) passing through \(x\) . Show that \(g\) is a continuous mapping from \(B\) onto \(S\) , such that \(g(x) = x\) for all \(x \in S\) (compare with GT, VI, § 2, exerc. 8). Deduce that there exists \(x_0 \in S\) and a continuous mapping \(h\) from \(S \times [0,1]\) onto \(S\) such that \(h(x,0) = x_0\) and \(h(x,1) = x\) for all \(x \in S\) .
\end{enumerate}

\begin{enumerate}
\def\labelenumi{\arabic{enumi})}
\setcounter{enumi}{13}
\tightlist
\item
  \begin{enumerate}
  \def\labelenumii{\alph{enumii})}
  \tightlist
  \item
    Let \((\mathrm{E}_n)_{n\geqslant 0}\) be an infinite sequence of real Banach spaces, E the vector subspace of the product \(\mathrm{F} = \prod_{n = 0}^{\infty}\mathrm{E}_n\) consisting of all sequences \(x = (x_{n})\) such that \(\sum_{n}\| x_{n}\|^{2} < + \infty\) . Show that the function \(\| x\| = (\sum_{n}\| x_{n}\|^{2})^{1 / 2}\) on E, is a norm, and that E is complete for this
  \end{enumerate}
\end{enumerate}

norm : we say that E is the hilbertian sum of the Banach spaces \(E_{n}\) .

\begin{enumerate}
\def\labelenumi{\alph{enumi})}
\setcounter{enumi}{1}
\item
  Show that the strong dual \(\mathbf{E}'\) of \(\mathbf{E}\) can be identified with the hilbertian sum of the strong duals \(\mathrm{E}_n'\) of the spaces \(\mathrm{E}_n\) , and that if \(x' = (x_n') \in \mathrm{E}'\) , then \(\langle x, x' \rangle = \sum_{n} \langle x_n, x_n' \rangle\) (if \(u\) is a continuous linear form on E, \(u_{n}\) its restriction to \(E_{n}\) considered as a subspace of E, and \(a_{n}\) a point of \(E_{n}\) such that \(\|a_{n}\|=1\) , show that, for every sequence \((\lambda_{n})\) of real numbers such that \(\sum_{n}\lambda_{n}^{2}<+\infty\) the series with the general term \(\lambda_{n}u_{n}(a_{n})\) is convergent, and deduce that \(\sum_{n}(u_{n}(a_{n}))^{2}<+\infty\) , using Banach-Steinhaus theorem for \(\ell^{2}(\mathbf{N})\) , for example).
\item
  Deduce from \(b)\) that when each of the \(E_n\) is reflexive, then \(E\) is reflexive. In particular, if for \(E_n\) we take the space \(\mathbf{R}^n\) endowed with the norm \(\| x \| = \sup_{1 \leqslant i \leqslant n} |\xi_i|\) for \(x = (\xi_i)_{1 \leqslant i \leqslant n}\) , show that E is reflexive, but there does not exist any norm compatible with the topology of E and for which E is uniformly convex (V, p.~67, exerc. 31).
\end{enumerate}

¶ 15) a) For every integer n \textgreater{} 0, let \(a^{(n)}\) be the double sequence defined in IV, p.~63, exerc. 8. Let E be the vector space of double sequences \(x = (x_{ij})\) of real numbers such that, for every integer n \textgreater{} 0, we have \(p_n(x) = (\sum_{i,j} a_{ij}^{(n)} |x_{ij}|^2)^{1/2} < +\infty\) . Show that the \(p_n\) are semi-norms on E,

and that E endowed with the topology defined by these semi-norms is a Fréchet space and a Montel space (argue as in IV, p.~60, exerc. 11).

\begin{enumerate}
\def\labelenumi{\alph{enumi})}
\setcounter{enumi}{1}
\item
  Show that the dual of E can be identified with the space \(E'\) of all double sequences \(x' = (x'_{ij})\) such that, for at least one index n, we have \(\sum_{i,j} (a_{ij}^{(n)})^{-1} |x'_{ij}|^2 < +\infty\) .
\item
  For every \(x = (x_{ij}) \in \mathbf{E}\) , show that \(\sum_{j=1}^{\infty} \left( \sum_{i=1}^{\infty} |x_{ij}|^2 \right) < +\infty\) (use Cauchy-Schwarz inequality); for every \(j \geqslant 1\) , put \(y_j = \sum_{i=1}^{\infty} x_{ij}\) ; the sequence \(u(x) = (y_j)\) then belongs to the hilbertian space \(\ell^2(\mathbf{N})\) . Show that \(u\) is a strict surjective morphism from \(E\) onto \(\ell^2(\mathbf{N})\) ; deduce that there exists weakly compact sets in \(\ell^2(\mathbf{N})\) which are not the images under \(u\) of a bounded set in \(E\) (argue as in IV, p.~63, exerc. 8).
\end{enumerate}

¶ 16) a) Let \(\Lambda\) be the set of increasing mappings \(\lambda : \mathbf{N} \to \mathbf{R}_{+}^{*}\) ; for every integer \(n \geqslant 0\) , and every \(\lambda \in \Lambda\) , let \(\phi_{n}(\lambda) = \lambda(n)\) . Let E be the set of all mappings \(x: \Lambda \to \mathbf{C}\) such that, for every \(n \in \mathbf{N}\) , we have \(p_{n}(x) = \left( \sum_{\lambda \in \Lambda} |x(\lambda)|^{2} \phi_{n}(\lambda) \right)^{1/2} < +\infty\) . Show that E is a vector space on which

the \(p_{n}\) are the semi-norms defining a reflexive Fréchet space structure.

\begin{enumerate}
\def\labelenumi{\alph{enumi})}
\setcounter{enumi}{1}
\item
  Let \(\mathbf{B}\) be a bounded set in \(\mathbf{E}\) , and let \(\alpha_{n} = \sup_{x\in \mathbf{B}}p_{n}(x)\) ; let \(\lambda_0\) be an element of \(\Lambda\) such that \(\lim_{n\to \infty}\lambda_0(n)^{-1}\alpha_n^2 = 0\) . Show that \(x(\lambda_0) = 0\) for all \(x\in \mathbf{B}\) , and hence that the set \(\mathbf{B}\) is not total in \(\mathbf{E}\) .
\item
  Let \((\mathrm{U}_n)\) be a countable fundamental system of convex and balanced neighbourhoods of 0 in E; if \(\mathrm{U}_n^\circ\) is metrizable for the strong topology on \(\mathrm{E}'\) , then there exists a sequence \((\mathrm{B}_{nm})_{m\geqslant 0}\) of bounded sets in E such that the sets \(\mathbf{B}_{nm}^{\circ}\cap \mathbf{U}_n^\circ\) form a fundamental system of neighbourhoods of 0 in \(\mathbf{U}_n^\circ\) for the strong topology. Deduce from \(b\) ) that there exists an integer \(n\) such that \(\mathrm{U}_n^\circ\) is not metrizable for the strong topology (use exerc. 5 of III, p.~38).
\end{enumerate}

